% Discussion 4.1 --- reinforcing findings: social context and thermal comfort.
\begin{table}[htbp]
\centering
\caption{Space-type findings that reinforce prior intuition: social context and thermal comfort. Details are drawn from the fingerprint, lift, and residual analyses of Section~\ref{sec:res-fingerprint-association}.}
\label{tab:insights-social-thermal}
\footnotesize
\renewcommand{\arraystretch}{1.3}
\begin{tabularx}{\textwidth}{@{}>{\raggedright\arraybackslash\bfseries}p{0.30\textwidth} >{\raggedright\arraybackslash}X@{}}
\toprule
\normalfont\textbf{Key insight} & \textbf{Detail} \\
\midrule
Home is overwhelmingly solitary & About 84\% of home responses report being \emph{alone}, the most solitary setting measured. \\
Classrooms are the most social setting & Classrooms are the most group-oriented space, with a group lift near 2.1. \\
Activity is the strongest spatial marker & Group activity is the response most tightly associated with space type (Cram\'er's V $\approx 0.38$). \\
Offices show the lowest cooling demand & Only about 26\% of office responses want \emph{cooler}, consistent with heavy air-conditioning. \\
Outdoor spaces show the highest cooling demand & Just over half ($\approx 53\%$) of outdoor responses want \emph{cooler}. \\
\bottomrule
\end{tabularx}
\end{table}
